\section{Problem Set 4: Systems of Linear Equations}

\begin{enumerate}
\item \textbf{Problem 1}\par
Solve by substitution:
\[
\begin{cases}
x+y=7,\\
2x-y=2.
\end{cases}
\]

\item \textbf{Problem 2}\par
Write the augmented matrix for
\[
\begin{cases}
2x-y+3z=5,\\
x+4y-z=2,\\
3x+2z=7.
\end{cases}
\]

\item \textbf{Problem 3}\par
Solve by Gaussian elimination:
\[
\begin{cases}
x+2y=5,\\
3x-y=4.
\end{cases}
\]

\item \textbf{Problem 4}\par
Solve by Gaussian elimination:
\[
\begin{cases}
x+y+z=6,\\
2x-y+z=3,\\
x+3y-2z=1.
\end{cases}
\]

\item \textbf{Problem 5}\par
Reduce the augmented matrix
\[
\left(
\begin{array}{cc|c}
1&2&3\\
2&4&8
\end{array}
\right)
\]
to echelon form and classify the system.

\item \textbf{Problem 6}\par
Reduce
\[
\left(
\begin{array}{ccc|c}
1&2&-1&3\\
2&4&-2&6\\
0&1&1&2
\end{array}
\right)
\]
and express all solutions using a parameter.

\item \textbf{Problem 7}\par
Determine whether
\[
\begin{cases}
x-2y=1,\\
2x-4y=5
\end{cases}
\]
has no solution, one solution, or infinitely many solutions. Interpret the result geometrically.

\item \textbf{Problem 8}\par
Solve
\[
\begin{cases}
2x+y=4,\\
x-3y=-5
\end{cases}
\]
using Cramer's rule.

\item \textbf{Problem 9}\par
Solve the system in the previous problem using the inverse matrix method. Compare the two methods.

\item \textbf{Problem 10}\par
For which values of \(k\) does
\[
\begin{cases}
kx+y=2,\\
2x+ky=3
\end{cases}
\]
have a unique solution?

\item \textbf{Problem 11}\par
Classify the system according to the value of \(a\):
\[
\begin{cases}
x+y=2,\\
2x+2y=a.
\end{cases}
\]

\item \textbf{Problem 12}\par
Find all values of \(k\) for which
\[
\begin{cases}
x+y+z=1,\\
2x+2y+2z=2,\\
x+y+kz=3
\end{cases}
\]
is consistent. Describe the solution set in each consistent case.

\item \textbf{Problem 13}\par
The row-reduced augmented matrix of a system is
\[
\left(
\begin{array}{cccc|c}
1&0&2&0&3\\
0&1&-1&0&4\\
0&0&0&1&2
\end{array}
\right).
\]
Identify the pivot and free variables and write the complete solution parametrically.

\item \textbf{Problem 14}\par
Determine whether the point \((2,-1)\) is the solution of
\[
\begin{cases}
3x+y=5,\\
x-2y=4.
\end{cases}
\]
If it is not, find the correct solution.

\item \textbf{Problem 15}\par
Find the intersection point of the lines
\[
2x-y-1=0,
\qquad
x+3y-9=0.
\]

\item \textbf{Problem 16}\par
Find the intersection of the three planes
\[
\begin{cases}
x+y+z=3,\\
x-y+z=1,\\
2x+z=2.
\end{cases}
\]

\item \textbf{Problem 17}\par
A cinema sells adult tickets for \(12\) units and student tickets for \(8\) units. In one evening it sells \(180\) tickets for a total of \(1840\) units. Form and solve a linear system to find the number of each type of ticket.

\item \textbf{Problem 18}\par
A system reduces to a row
\[
\begin{pmatrix}0&0&0\mid 6\end{pmatrix}.
\]
What does this row mean, and what can be concluded about the original system?

\item \textbf{Problem 19}\par
Construct a system of two linear equations in two variables that has:
\begin{enumerate}
\item exactly one solution;
\item no solution;
\item infinitely many solutions.
\end{enumerate}
Explain how the coefficients produce each case.

\item \textbf{Problem 20}\par
Consider
\[
\begin{cases}
x+ay=1,\\
ax+y=1.
\end{cases}
\]
Classify the system for every real value of \(a\). For each case, determine whether the solution is unique, nonexistent, or non-unique, and find the solution set where possible.
\end{enumerate}
